\chapter{Data types}

\section{Core}

\section{Basic}
    
    \subsection{Bytes}
        This is the rawest form of data
        \begin{lstlisting}[language=Python]
        raw = b'\x48\x65\x6c\x6c\x6f'
        raw.decode("utf-8")  # "Hello"
        data_bytes = b'{"name": "Ali", "age": 30}'
        data_dict = json.loads(data_bytes.decode("utf-8"))
        print(data_dict["name"])  # Ali
        \end{lstlisting}
        
        \paragraph{As data type hint}
            \begin{lstlisting}[language=Python]
                class Interface(Protocol):
                    _ram: bytes
            \end{lstlisting}
        
        \paragraph{As data type}
            \begin{lstlisting}[language=Python]
                data_bytes = b'hello'
                data_str = data_bytes.decode("utf-8")
                print(data_str)   # "hello"
            \end{lstlisting}
        
        \paragraph{Converting to string}
            \begin{lstlisting}[language=Python]
                data_str = "hello"
                data_bytes = data_str.encode("utf-8")
                print(data_bytes)  # b'hello'
            \end{lstlisting}
    
    \subsection{Numbers}
        
        \paragraph{int}
        
        \paragraph{float}
        
        \paragraph{complex}
        
        \paragraph{Decimal}
            for arbitary precision fro scientific calculation
        
        \paragraph{Fraction}
            For rational numbers

\section{Mutability}

\section{Collections}
    \subsection{Sequence}
        They have order
    
        \paragraph{list}
        \paragraph{tuple}
        \paragraph{str}
        \paragraph{range}
    \subsection{Shared functions}
        \paragraph{slice}
            \begin{minted}{python}
                s = slice(2, 5, 1)
                nums = [0, 1, 2, 3, 4, 5, 6, 7, 8, 9]
                print(nums[s])
            \end{minted}
    \subsection{set types}
        \paragraph{set}
        \paragraph{fozenset}
    
    \subsection{Mapping types}
        Mapping type is a key value type
        \paragraph{dict}
        \paragraph{defaultdict}
            If the key doesnt exist it creates it
        \paragraph{ChainMap}
            Combines two dicts
        \paragraph{MappingProxyType}
            A read only dict
        \paragraph{Ordered Dictoionary}
            \begin{minted}{python}
                from collections import OrderedDict
    
                od = OrderedDict()
                od["a"] = 1
                od["b"] = 2
                od["c"] = 3
                
                print(od)  # OrderedDict([('a', 1), ('b', 2), ('c', 3)])
    
                
                od.move_to_end("b")
                print(od)  # OrderedDict([('a', 1), ('c', 3), ('b', 2)])
    
                od.move_to_end("b", last=False)
                print(od)  # OrderedDict([('b', 2), ('a', 1), ('c', 3)])
            \end{minted}
    

\section{Hashability}
    A hashable object has implemented \pyl{__hash__()} and \pyl{__eq__()} functions. Data types such as int, str, float are hashable but list, dict and set are not
    \paragraph{\_\_hash\_\_()}
        when implemented in a class, then it should return an  integer value sokmehow related to that object

\section{methods}
    
    remove(element): raises key error if element doesnt exist
    
    discard(element): removes but does nothing if the element doesnt exist
    
    pop: removes and returns an arbitrary element
    
    clear: removes all the elements

\section{type hints}
    \paragraph{Annotated}
        with this type you can put comments for example in the following example x muist be positive
        \begin{minted}{python}
            from typing import Annotated
            x: Annotated[int, "positive"] = 42
        \end{minted}
        with fast API, this example say that the integer myst be between 0 and 100
        \begin{minted}{python}
            from typing import Annotated
            from fastapi import Query
            
            def get_items(limit: Annotated[int, Query(gt=0, lt=100)]):
                return {"limit": limit}
        \end{minted}

    \paragraph{Optional}
        
        When the type is a class object or None
        
        \begin{lstlisting}[language=Python]
            ->Optional[CommentedMap]
        \end{lstlisting}
        
        The above code says the return is either CommentedMap or None

    \paragraph{Generic}
        \begin{lstlisting}[language=Python]
            from typing import Generic, TypeVar, Protocol
            T = TypeVar("T")
            class Unit(Protocol, Generic[T]):
                def __iter__(self) -> Iterator[T]: ...
                def __len__(self) -> int: ...
        \end{lstlisting}